% ===================================================================
%  TABLEAU N° 2 — Le Corps humain. — La Récréation. — Les Jeux.
%  Feuillets 14 a 18 du fac-simile = folios 12 a 16.
%  Correspondance : folio imprime = numero de feuillet - 2.
%
%  Les cles « %%K » sont les MEMES que dans texto/io : c'est par elles
%  que la page de lecture met les deux textes en regard. Le francais ne
%  connait pas la subdivision idote en « La Kapo », « La Torso », « La
%  Membri », ni les sous-titres « Gimnastiko », « La Ludi en liber
%  aero », « Biciklago e balonludo » : la table francaise se limite a
%  trois intertitres (I. Le Corps humain, II. La Gymnastique, III. Les
%  Jeux), et ces cles-la restent donc sans vis-a-vis. C'est un ecart
%  REEL entre les deux editions, non une lacune du releve.
%
%  ATTENTION — LE FAC-SIMILE FRANCAIS EST UNE PRISE DE VUE, non un
%  passage au scanner : contraste inegal, page gondolee. Chaque page a
%  ete verifiee a l'agrandissement (2400 px) avant releve.
% ===================================================================

% --- Feuillet 14 (folio 12 : ouverture de tableau) ------------------
\begin{VUpage}[14]{12}
%%K t02-apar-1 apar
\VUsaut{2.0mm}
\VUtitre{15.0pt}{60}{TABLEAU N\textsuperscript{o} 2}
\VUsaut{2.4mm}
\VUtitre{11.4pt}{0}{Le Corps humain. --- La Récréation. --- Les Jeux}
\VUsaut{4.0mm}

%%K t02-tit-1 sub
\VUpk{10.2pt}{40}{I. Le Corps Humain.}
\VUsaut{2.4mm}
\VUcentre{10.2pt}{0}{\textit{(Simple leçon d'histoire naturelle).}}
\VUsaut{3.4mm}

%%K t02-01-1 p
\VUblancAlinea
1. --- Les principales parties du corps humain sont:\nl
la \VUgras{tête,} le \VUgras{tronc} et les \VUgras{membres.}

%%K t02-02-1 p
\VUblancAlinea
2. --- La tête comprend deux parties : le \VUgras{crâne,}\nl
qui renferme le \VUgras{cerveau,} et que les \VUgras{cheveux} recou\cc
vrent, et le \VUgras{visage} (la figure).

%%K t02-03-1 p
\VUblancAlinea
3. --- Sur notre dessin nous ne voyons ni le crâne\nl
ni le cerveau ; mais nous apercevons très nettement\nl
les parties importantes du visage.

%%K t02-04-1 p
\VUblancAlinea
4 --- Voici d'abord le \VUgras{front,} qui annonce une\nl
intelligence puissante, une pensée toujours active.\nl
Les \VUgras{tempes} sont très dégagées. L'\VUgras{œil} est vif et\nl
bien ouvert. Un épais \VUgras{sourcil} et de long \VUgras{cils} le\nl
protègent.

%%K t02-05-1 p
\VUblancAlinea
5. --- Voici encore le \VUgras{nez} et les \VUgras{narines.} Le nez\nl
nous sert à sentir et à respirer. De chaque côté du\nl
nez sont les \VUgras{joues} ; plus loin, les \VUgras{oreilles.} Le bas\nl
de la figure est formé par la \VUgras{bouche} et le \VUgras{menton.}

%%K t02-06-1 p
\VUblancAlinea
6. --- Nous ne les voyons pas très bien, car la\nl
\VUgras{moustache} et la \VUgras{barbe} nous les cachent. Mais\nl
nous savons que la bouche comprend les deux\nl
\VUgras{lèvres,} les \VUgras{mâchoires,} qui portent les \VUgras{dents,}\nl
et la \VUgras{langue.} Avec la bouche nous parlons et nous\nl
mangeons. Avec la langue nous goûtons les aliments.

%%K t02-07-1 p
\VUblancAlinea
7. --- L'œil, l'oreille, le nez et la bouche sont, avec\nl
les doigts, les organes des cinq sens : la \VUgras{vue,} l'\VUgras{ouïe,}\nl
l'\VUgras{odorat,} le \VUgras{goût} et le \VUgras{toucher.}

%%K t02-08-1 p
\VUblancAlinea
8. --- La tête est donc une des parties les plus\nl
intéressantes du corps humain.
\end{VUpage}

% --- Feuillet 15 (folio 13) -----------------------------------------
\begin{VUpage}[15]{13}
%%K t02-09-1 p
\VUblancAlinea
9. --- Le \VUgras{cou} réunit la tête et le tronc. Il comprend\nl
la \VUgras{gorge} et la \VUgras{nuque.}

%%K t02-10-1 p
\VUblancAlinea
10. --- Le tronc contient aussi beaucoup d'organes\nl
essentiels. Dans la \VUgras{poitrine} sont les \VUgras{poumons,} avec\nl
lesquels nous respirons, et le \VUgras{cœur,} qui est le centre\nl
de la vie. Lorsque le cœur cesse de battre, l'être\nl
vivant meurt.

%%K t02-11-1 p
\VUblancAlinea
11. --- Les \VUgras{côtes} supportent la poitrine.

%%K t02-12-1 p
\VUblancAlinea
12. --- Les autres parties du tronc sont le \VUgras{flanc}\nl
(côté), le \VUgras{dos,} les \VUgras{épaules} et le \VUgras{ventre} (abdomen).\nl
Nous nous couchons sur le côté ou sur le dos pour\nl
dormir, nous portons les fardeaux sur notre épaule.

%%K t02-13-1 p
\VUblancAlinea
13. --- Dans l'abdomen se trouvent l'\VUgras{estomac} et\nl
les \VUgras{intestins.} Il nous servent à digérer les aliments.

%%K t02-14-1 p
\VUblancAlinea
14. --- Nous avons quatre \VUgras{membres,} deux \VUgras{bras}\nl
et deux \VUgras{jambes.} Avec les bras nous prenons, nous\nl
tenons, nous soulevons, nous ramassons les objets;\nl
avec les jambes nous nous tenons debout, nous mar\cc
chons et nous courons.

%%K t02-15-1 p
\VUblancAlinea
15. --- L'\VUgras{épaule} unit le bras au corps. Un \VUgras{muscle}\nl
très fort, le biceps, fait mouvoir le bras que le \VUgras{coude}\nl
relie à l'\VUgras{avant-bras.} Le \VUgras{poignet} constitue l'articu\cc
lation de la \VUgras{main} qui se termine par les \VUgras{doigts} et\nl
les \VUgras{ongles.} Sur la main nous distinguons une \VUgras{veine}\nl
(artère) dans laquelle coule le \VUgras{sang.}

%%K t02-16-1 p
\VUblancAlinea
16. --- Les parties de la jambe sont : la \VUgras{cuisse,} le\nl
\VUgras{genou} et le \VUgras{jarret,} le \VUgras{mollet} dont la \VUgras{chair} est\nl
recouverte par la \VUgras{peau.} la \VUgras{cheville} et le \VUgras{pied.} Les\nl
\VUgras{os} de la jambe sont très gros et très forts. A l'ex\cc
trémité du pied sont les \VUgras{doigts de pied.} Les autres\nl
parties sont la \VUgras{plante du pied} et le \VUgras{talon.}

%%K t02-17-1 p
\VUblancAlinea
17. --- Telle est la description à peu près complète\nl
du corps humain.

%%K t02-17-2 p
\VUblancAlinea
\textit{Nota.} --- \textit{A l'aide de l'expression : J'ai mal à...}\nl
\textit{(la tête, etc.), le professeur pourra reprendre tout}\nl
\textit{ce vocabulaire au point de vue de la} \VUgras{santé.}
\end{VUpage}

% --- Feuillet 16 (folio 14) -----------------------------------------
\begin{VUpage}[16]{14}
%%K t02-tit-2 sub
\VUpk{10.2pt}{40}{II. La Gymnastique.}
\VUsaut{2.4mm}
\VUcentre{10.2pt}{0}{\textit{(Récit d'un élève).}}
\VUsaut{3.4mm}

%%K t02-18-1 p
\VUblancAlinea
18. --- Pour être souples, agiles et bien portants,\nl
nous devons exercer nos jambes et nos bras. C'est\nl
pourquoi nous jouons et nous faisons de la \VUgras{gymnas}\cc
\VUgras{tique.}

%%K t02-19-1 p
\VUblancAlinea
19. --- En semaine, ces deux exercices ont lieu\nl
dans la cour du Lycée (voir Tableau n\textsuperscript{o} 1). Mais, le\nl
jeudi et le dimanche, nous allons sur une grande\nl
\textit{pelouse,} où nous avons beaucoup de place pour courir\nl
et nous amuser.

%%K t02-20-1 p
\VUblancAlinea
20. --- Nos maîtres et nos parents nous y accompa\cc
gnent parfois ; mais c'est le \VUgras{professeur de gym}\cc
\VUgras{nastique} \textsuperscript{(12)} qui dirige les exercices.

%%K t02-21-1 p
\VUblancAlinea
21. --- Sur la pelouse, à gauche en entrant, se\nl
trouvent les \VUgras{appareils de gymnastique} \textsuperscript{(1)}, nous\nl
grimpons à l'\VUgras{échelle} \textsuperscript{(2)}, nous faisons des \VUgras{renverse}\cc
\VUgras{ments} aux \VUgras{anneaux} \textsuperscript{(3)} et au \VUgras{trapèze} \textsuperscript{(4)} ; nous nous\nl
hissons avec les mains jusqu'au sommet de \VUgras{l'échelle}\nl
\VUgras{de corde} \textsuperscript{(5)} ou de la \VUgras{corde à nœuds} \textsuperscript{(6)}.

%%K t02-22-1 p
\VUblancAlinea
22. --- Pendant ce temps nos camarades sautent\nl
par dessus le \VUgras{cheval de bois} \textsuperscript{(7)} ou les \VUgras{barres}\nl
\VUgras{parallèles} \textsuperscript{(11)}. D'autres font de la \VUgras{boxe} \textsuperscript{(8)} ou de\nl
l'\VUgras{escrime} \textsuperscript{(9)} avec un masque et un \VUgras{fleuret.} D'autres\nl
enfin soulèvent des \VUgras{haltères} \textsuperscript{(10)}.

%%K t02-tit-3 sub
\VUsaut{5.0mm}
\VUpk{10.2pt}{40}{III. --- Les Jeux.}
\VUsaut{3.4mm}

%%K t02-23-1 p
\VUblancAlinea
23. --- Autour des gymnastes, garçons et fillettes\nl
jouent à différents \VUgras{jeux.} Les petites filles s'amusent\nl
avec leur \VUgras{cerceau} \textsuperscript{(14)} ; elles sautent à la \VUgras{corde} \textsuperscript{(13)},\nl
ou bien elles invitent leurs frères et leurs cousins\nl
à faire une partie de \VUgras{croquet} \textsuperscript{(15)}. Elles sont enchan\cc
tées de croquer la \VUgras{boule} de l'adversaire avec leur\nl
\VUgras{maillet} au moment où il croit traverser facilement\nl
un \VUgras{arceau} !
\end{VUpage}

% --- Feuillet 17 (folio 15) ------------------------------------------
\begin{VUpage}[17]{15}
%%K t02-24-1 p
\VUblancAlinea
24. --- Les garçons aiment mieux les jeux plus\nl
violents. Ils jouent volontiers aux \VUgras{quilles} \textsuperscript{(16)}, qu'ils\nl
abattent adroitement avec la \VUgras{boule} \textsuperscript{(17)}. Ils ne détes\cc
tent pas non plus la \VUgras{toupie} \textsuperscript{(18)} et le \VUgras{bilboquet} \textsuperscript{(19)}\nl
ou même le \VUgras{jeu de tonneau} \textsuperscript{(20)}. Mais leurs jeux\nl
préférés sont encore les \VUgras{billes} \textsuperscript{(21)} et la \VUgras{balle au}\nl
\VUgras{mur} \textsuperscript{(22)}, car le \VUgras{mur} de la \VUgras{serre} \textsuperscript{(26)} est très commode\nl
pour faire rebondir la balle légère.

%%K t02-25-1 p
\VUblancAlinea
25. --- Le petit Jean, perché sur ses \VUgras{échasses} \textsuperscript{(24)},\nl
est très adroit et sait fort bien garder l'équilibre.\nl
Près de lui, quelques-uns de ses camarades jouent à\nl
\VUgras{cache-cache} \textsuperscript{(25)} dans la serre et derrière la \VUgras{haie.}\nl
D'autres préfèrent la \VUgras{main-chaude} \textsuperscript{(28)}, ou le \VUgras{vo}\cc
\VUgras{lant} \textsuperscript{(29)}, qui voltige légèrement d'une \VUgras{raquette} à\nl
l'autre.

%%K t02-26-1 p
\VUblancAlinea
26. --- Au milieu de la pelouse il y a deux arbres.\nl
Au pied de l'un d'eux, sept ou huit garçons ont\nl
organisé une partie de \VUgras{cheval fondu} \textsuperscript{(30)}. On ne\nl
connaît pas ce jeu dans tous les pays ; mais tout le\nl
monde sait ce qu'est le jeu de \VUgras{saute-mouton,} auquel\nl
les enfants ne jouent pas en ce moment.

%%K t02-27-1 p
\VUblancAlinea
27. --- Le \VUgras{colin-maillard} \textsuperscript{(31)} est aussi fort amu\cc
sant ; fillettes et garçons mettent tour à tour le\nl
\VUgras{bandeau} sur leurs yeux pour attraper les maladroits\nl
qui se laissent prendre.

%%K t02-28-1 p
\VUblancAlinea
28: --- Au milieu de la pelouse, voici un jeu bien\nl
français, mais un peu brutal : c'est \VUgras{l'ourse} \textsuperscript{(32)}. Il est\nl
beaucoup plus bruyant que le jeu si paisible du \VUgras{cerf}\cc
\VUgras{volant} \textsuperscript{(33)} ou même que le jeu anglais du \VUgras{lawn}\cc
\VUgras{tennis} \textsuperscript{(34)}.

%%K t02-29-1 p
\VUblancAlinea
29. --- Ce dernier sport fait la joie des grands\nl
élèves. Nos professeurs eux-mêmes s'y livrent volon\cc
tiers. Il demande beaucoup d'adresse et d'agilité et me\nl
plaît fort. Mais je laisse aux petites filles les plaisirs\nl
de la \VUgras{balançoire} \textsuperscript{(35)}.

%%K t02-30-1 p
\VUblancAlinea
30. --- Je n'aime pas beaucoup non plus un autre\nl
jeu anglais qui peut devenir dangereux. Je veux
\parplein
\end{VUpage}

% --- Feuillet 18 (folio 16) --------------------------------------------
\begin{VUpage}[18]{16}
%%K t02-30-1b p
\VUcontinue
parler du \VUgras{foot-ball} \textsuperscript{(36)}, qui fait le bonheur de\nl
mon frère aîné. J'aime mieux notre vieux \VUgras{jeu de}\nl
\VUgras{barres} \textsuperscript{(38)}, tout aussi hygiénique et bien moins\nl
brutal.

%%K t02-31-1 p
\VUblancAlinea
31. --- Je fais aussi souvent un tour de pelouse à\nl
\VUgras{bicyclette} \textsuperscript{(37)}.

%%K t02-32-1 p
\VUblancAlinea
32. --- L'année prochaine j'apprendrai \VUgras{l'équita}\cc
\VUgras{tion} \textsuperscript{(39)}. C'est si beau un cheval qui marche gracieu\cc
sement au pas ou au trot, ou bien qui saute les\nl
\VUgras{obstacles} au galop !

%%K t02-33-1 p
\VUblancAlinea
33. --- En été nous apprenons la \VUgras{natation} \textsuperscript{(40)}.\nl
Nous plongeons et nous nous baignons avec délices\nl
dans l'eau claire et fraîche de la petite rivière.\nl
Parfois enfin nous lançons un \VUgras{ballon} en papier \textsuperscript{(41)} qui\nl
s'élève majestueusement dans les airs dès qu'il est\nl
plein d'air chaud.

%%K t02-34-1 p
\VUblancAlinea
34. --- Il y a encore bien d'autres jeux. \VUgras{Mais je}\nl
n'en finirais pas de les nommer, et l'on voit par \VUgras{ce}\nl
simple aperçu que l'on ne s'ennuie guère le jeudi \VUgras{sur}\nl
notre belle pelouse.

%%K t02-34-2 p
\VUblancAlinea
\textit{N. B.} --- \textit{Le professeur complétera facilement ce}\nl
\textit{chapitre en décrivant plus minutieusement} \VUgras{\textit{chacun des}}\nl
\textit{jeux importants.}

\VUsaut{6.0mm}
\VUfilet{22mm}
\end{VUpage}
